\chapter{Allergies to Seafood}
\index{Allergies to Seafood}

\section*{Definition and Overview}
A seafood allergy is an immunoglobulin E (IgE)-mediated immune reaction to proteins present in fish or shellfish. The term covers two biologically distinct groups: finned fish, such as salmon, tuna, cod, and sardines, and shellfish, which is further divided into crustaceans (prawns, crab, and lobster) and molluscs (oysters, mussels, clams, and squid). Allergy to one group does not reliably predict allergy to the other, although significant cross-reactivity occurs within each group. Reactions typically begin within minutes to two hours of ingestion and range from mild oral itching and hives to severe anaphylaxis. Because seafood allergy is usually lifelong, management centres on strict avoidance, careful label reading, and preparedness for accidental exposure.

\section*{Epidemiology}
Seafood allergy is one of the most common food allergies in adults and is especially prevalent in regions with high seafood consumption, including parts of Asia, and Southern Europe. It is estimated to affect roughly 1 to 3\% of the general population in many developed countries. In contrast to allergies to cow's milk and egg, which are typically outgrown, seafood allergy commonly begins in adulthood and usually persists for life. Crustacean allergy is reported more often than fish or mollusc allergy, and people allergic to one species within a group frequently react to related species. Occupational exposure among fish-processing and catering workers is a recognised additional cause of sensitisation.

\section*{Aetiology and Pathophysiology}
Seafood allergy is an IgE-mediated reaction to specific proteins in the flesh of fish and shellfish:

\begin{itemize}
    \item \textbf{Fish proteins:} parvalbumin is the major allergen in most finned fish and is highly cross-reactive between species; other allergens include enolase and aldolase.
    \item \textbf{Crustacean proteins:} tropomyosin is the dominant allergen in prawns, crab, and lobster, and is also the basis for cross-reactivity with molluscs and with dust mite tropomyosin.
    \item \textbf{Mollusc proteins:} similar tropomyosin allergens occur in oysters, mussels, clams, and squid.
\end{itemize}

On exposure, the allergen cross-links specific IgE on mast cells, triggering release of histamine and other mediators within minutes to a couple of hours. In some people, symptoms can also be provoked by inhaling steam from cooking seafood or by handling it directly, which is particularly relevant for occupational allergy.

\section*{Clinical Features}
Symptoms usually appear soon after eating seafood:

\begin{itemize}
    \item \textbf{Skin:} urticaria, itching, and swelling of the lips, face, or tongue.
    \item \textbf{Mouth and throat:} tingling or itching of the mouth, palate, and throat.
    \item \textbf{Gut:} nausea, vomiting, abdominal cramps, and diarrhoea.
    \item \textbf{Respiratory:} runny nose, sneezing, red or watery eyes, cough, wheeze, and difficulty breathing.
    \item \textbf{Severe reactions:} throat swelling, dizziness, hypotension, and collapse, representing anaphylaxis.
\end{itemize}

The severity of a reaction varies between individuals and between episodes, and a person who has previously had only mild symptoms can still experience a severe reaction on another occasion.

\section*{Differential Diagnosis}
Not every illness after eating seafood is an allergy:

\begin{itemize}
    \item \textbf{Food poisoning:} from undercooked or spoiled fish and shellfish, causing vomiting and diarrhoea without allergic features.
    \item \textbf{Scombroid poisoning:} caused by histamine that accumulates in poorly stored fish such as tuna and mackerel, closely mimicking an allergic reaction.
    \item \textbf{Shellfish poisoning:} from naturally occurring marine toxins associated with algal blooms.
    \item \textbf{Viral gastroenteritis:} particularly norovirus, which is often linked to shellfish.
    \item \textbf{Anisakis infestation:} from raw or undercooked fish, causing abdominal pain and sometimes allergic-type symptoms.
    \item \textbf{Other allergens:} a reaction to sauces, spices, or additives used with the seafood rather than the seafood itself.
\end{itemize}

\section*{When to Seek Medical Care}
Seek medical review if you suspect a seafood allergy, especially after any reaction, so the diagnosis can be confirmed and an individualised management plan developed. Do not attempt trial-and-error reintroduction of seafood at home.

\textbf{Contact your local emergency services for:}
\begin{itemize}
    \item Difficulty breathing, wheezing, or a hoarse voice after eating seafood.
    \item Swelling of the tongue, face, or throat, or difficulty swallowing.
    \item Dizziness, collapse, or loss of consciousness.
    \item Any signs of anaphylaxis, even if they seem mild at first.
\end{itemize}

\section*{Diagnosis and Investigations}
A specialist assessment is important to confirm the diagnosis and distinguish allergy from other causes:

\begin{itemize}
    \item \textbf{Detailed history:} the type of seafood eaten, the timing and nature of symptoms, prior reactions, and any occupational exposure.
    \item \textbf{Skin prick testing:} with fish and shellfish extracts to detect sensitisation.
    \item \textbf{Serum specific IgE:} blood tests for antibodies against seafood proteins, including parvalbumin and tropomyosin where available.
    \item \textbf{Oral food challenge:} supervised ingestion in a controlled clinical setting when the diagnosis remains unclear from testing.
    \item \textbf{Referral:} to an allergy specialist or clinical immunologist for confirmation and ongoing management.
\end{itemize}

\section*{Management}
Because seafood allergy rarely resolves, management focuses on avoidance and preparedness:

\begin{itemize}
    \item \textbf{Strict avoidance} of the confirmed seafood type, including in sauces, soups, stocks, and condiments such as fish sauce.
    \item \textbf{Reading labels:} checking for fish, prawn, crab, lobster, and mollusc in packaged and takeaway foods, and rechecking regularly as recipes change.
    \item \textbf{Cross-contamination awareness:} in restaurants, fishmongers, and shared cooking surfaces, and in frying oil shared between dishes.
    \item \textbf{Carrying an adrenaline auto-injector} and an up-to-date allergy action plan if prescribed.
    \item \textbf{Informing restaurant staff} clearly when ordering, and choosing venues that handle food allergies carefully.
    \item \textbf{Avoiding handling seafood} or being exposed to cooking steam if this triggers symptoms.
\end{itemize}

\section*{Complications}
Potential complications include:

\begin{itemize}
    \item Anaphylaxis from accidental exposure, particularly when eating out.
    \item Reactions from hidden ingredients or cross-contamination.
    \item Anxiety about eating out, travelling, or attending social events.
    \item Nutritional concerns if seafood is removed from the diet without substitution of alternative protein and omega-3 sources.
    \item Difficulty ensuring safe food when eating at other people's homes.
\end{itemize}

\section*{Prognosis and Outcomes}
Seafood allergy usually persists for life, in contrast to many childhood food allergies. It is nevertheless a very manageable condition for most people, because reactions occur only after exposure to seafood and avoidance is generally straightforward. Daily life is usually unaffected, provided there is a reliable action plan and an accessible adrenaline auto-injector where indicated. Regular review with an allergy specialist is recommended, and repeat testing may be offered to monitor whether the allergy is changing, although spontaneous resolution is uncommon.

\section*{Prevention and Self-Care}
Self-care centres on safe eating and being prepared for an accident:

\begin{itemize}
    \item Carry your adrenaline auto-injector and action plan at all times if prescribed.
    \item Wear medical alert identification describing your seafood allergy.
    \item Ask detailed questions when eating out and choose venues that take allergies seriously.
    \item Check labels on every occasion, since recipes and manufacturing processes change.
    \item Carry safe snacks when travelling or attending events.
    \item Keep an up-to-date written plan agreed with your specialist.
\end{itemize}

\section*{Sources and Further Reading}
The following reputable sources provide reliable, up-to-date information on seafood allergy:

\begin{itemize}
    \item NHS. \textit{Food allergy information, including fish and shellfish allergy.} https://www.nhs.uk/
    \item Mayo Clinic. \textit{Food allergy: symptoms, causes, and treatment.} https://www.mayoclinic.org/
    \item MSD Manual. \textit{Food allergy overview.} https://www.msdmanuals.com/
    \item MedlinePlus. \textit{Shellfish and fish allergy health information.} https://medlineplus.gov/
    \item NICE. \textit{Clinical guidance on food allergy in children and young people.} https://www.nice.org.uk/
    \item World Health Organization. \textit{Information on food safety and foodborne illness.} https://www.who.int/
\end{itemize}
